\section{A geometric bound at termination}\label{sec:accounting}

We now place the recursive state inside the ordinary unit cube $(0,1)^n$,
with one coordinate for each original label. The map preserves every
represented order. When sorting is complete its image must lie in a single
coordinate-order region, of volume $1/n!$. Computing the volume of the image
will therefore give an upper bound on the final state.

\subsection{An interval for each child}

Reserve for an anchor of address length $m$ a child interval of width
\begin{equation}\label{eq:child-weight}
w(m)=\frac{2^{-(1+1/k)m}}{4k}.
\end{equation}
At most $2^m$ anchors have that address length. Hence at every node
\begin{equation}\label{eq:weight-sum}
\sum_c w_c\le\frac1{4k}\sum_{m\ge0}2^{-m/k}
=\frac1{4k(1-2^{-1/k})}\le\frac12.
\end{equation}
For the last inequality, convexity gives $2^{-t}\le1-t/2$ for
$0\le t\le1$, and we take $t=1/k$.

At a node, map a free position $t$ distinct from all anchors to
\begin{equation}\label{eq:packing-map}
G(t)=\frac t2+\sum_{c<t}w_c.
\end{equation}
The jump at an anchor $c$ leaves the open interval
\[
J_c=\left(\frac c2+\sum_{a<c}w_a,
          \frac c2+\sum_{a<c}w_a+w_c\right).
\]
Place the recursively mapped child state inside $J_c$, using an affine map
from $(0,1)$ to $J_c$. The intervals $J_c$ are disjoint, occur in anchor
order, and fit in $(0,1)$ by \eqref{eq:weight-sum}. Free positions stay in
the gaps between them. All these widths are used only in the proof.

\begin{figure}[tb]
\centering
\begin{tikzpicture}[x=11cm,y=1cm,>=Latex,font=\small]
\draw[thick] (0,0)--(1,0);
\node[anchor=east] at (-.03,0) {parent};
\foreach \p/\s in {0/0,.28/{c_1},.7/{c_2},1/1}{
  \draw (\p,.05)--(\p,-.05); \node[above] at (\p,.1) {$\s$};}
\fill (.28,0) circle(1.6pt); \fill (.7,0) circle(1.6pt);
\draw[thick] (0,-1.25)--(1,-1.25);
\node[anchor=east] at (-.03,-1.25) {image};
\filldraw[draw=accent,fill=pale] (.14,-1.45) rectangle (.32,-1.05);
\filldraw[draw=accent,fill=pale] (.53,-1.45) rectangle (.67,-1.05);
\node at (.23,-1.25) {$J_{c_1}$};
\node at (.60,-1.25) {$J_{c_2}$};
\draw[->,accent] (.28,-.12)--(.23,-.94);
\draw[->,accent] (.70,-.12)--(.60,-.94);
\draw[decorate,decoration={brace,mirror,amplitude=4pt}]
  (.14,-1.55)--(.32,-1.55) node[midway,below=5pt] {$w_{c_1}$};
\draw[decorate,decoration={brace,mirror,amplitude=4pt}]
  (.53,-1.55)--(.67,-1.55) node[midway,below=5pt] {$w_{c_2}$};
\node[below] at (0,-1.30) {$0$};
\node[below] at (1,-1.30) {$1$};
\end{tikzpicture}
\caption{The packing map has slope $1/2$ between anchors. Its jump at each
anchor reserves an interval for the entire child subtree. The widths are
shown schematically; their values are given by \eqref{eq:child-weight}.}
\label{fig:packing}
\end{figure}

For a parent--child edge $c$, let $s_c$ be the number of original labels in
the child's entire subtree, and put
\begin{equation}\label{eq:W}
\mathcal W=\prod_{\text{parent--child edges }c}w_c^{s_c}.
\end{equation}
A comparison does not change $\mathcal W$. A refinement into a child of
width $w$ multiplies it by $w$: the refined label crosses that edge, while
its membership in every other subtree remains unchanged. This holds for
both a new child and an existing child.

\Needspace{12\baselineskip}
\begin{lemma}[Packing volume]\label{lem:packing}
The recursive map is injective and preserves the represented order. Its
image $\mathcal I\subseteq(0,1)^n$ has volume
\begin{equation}\label{eq:packed-volume}
\vol(\mathcal I)=2^{-n}Z\mathcal W.
\end{equation}
Consequently, if the transcript determines a unique total order, then
\begin{equation}\label{eq:terminal}
n!Z\mathcal W\le2^n.
\end{equation}
\end{lemma}
\begin{proof}
The map $G$ puts each free position on the correct side of every child
interval. Recursion therefore preserves the order of the original labels.
Anchors and widths are fixed within the state. The image coordinate of a
label depends only on its one free coordinate, through a strictly increasing
map, so the whole map is injective.

Partition the domain at the finitely many anchor hyperplanes. On each piece,
the map of a label is affine. Its slope is $1/2$ at its free coordinate,
multiplied by the widths of the child intervals containing that label along
the path to the root. There are exactly $n$ factors of $1/2$, one per free
coordinate. An edge width $w_c$ occurs $s_c$ times. Thus the Jacobian on
every piece is $2^{-n}\mathcal W$. The domain, a product of the node regions,
has volume $Z$. Injectivity lets us add the image volumes of the pieces,
proving \eqref{eq:packed-volume}.

At completion, legality places the image in the region of the cube realizing
the unique transcript order. The $n!$ coordinate-order regions have equal
volume $1/n!$, since coordinate permutations preserve volume and equality
hyperplanes have volume zero. Hence $2^{-n}Z\mathcal W\le1/n!$.
\end{proof}

\subsection{The addresses of a path}

The same bound on child widths controls a path repair.

\begin{lemma}[Address length of a path batch]\label{lem:path-address}
The $H$ labels of a selected directed path enter distinct children of their
parent. If their address lengths are $m_1,\ldots,m_H$, then
\begin{equation}\label{eq:path-charge}
\sum_{i=1}^H m_i\ge kH.
\end{equation}
\end{lemma}
\begin{proof}
Every arrow on the selected path is an answered comparison. Its labels are
therefore pairwise comparable in the public transcript, and remain so after
arrows are deleted. By Lemma~\ref{lem:admission}, two of these labels cannot
be admitted to the same child. The procedure completes all $H$ refinements
in this parent, so their distinct child widths have sum at most $1/2$.
The arithmetic--geometric mean inequality gives
\[
\prod_{i=1}^H w(m_i)\le(1/(2H))^H.
\]
Substitute \eqref{eq:child-weight} and take logarithms:
\[
(1+1/k)\sum_i m_i\ge H\bigl(\log H-\log(2k)\bigr).
\]
Since $H=2^{2k}$ and $\log(2k)\le k-1$ for integers $k\ge4$, the
right-hand side is at least $(k+1)H$. Division by $1+1/k$ proves the claim.
The auxiliary inequality follows from $2k\le2^{k-1}$, true at $k=4$ and
preserved when $k$ increases by one.
\end{proof}
